\scanpage{140}
% archive sheet number: 139

Il suffit de vérifier que l'on a bien

\[
tt'=b^2c'+b'{}^2c-4cc'\qquad\text{(\unsure{vu|on} qu'on a déjà }t+t'=bb'\text{)}
\]

Vérification immédiate (il n'est pas besoin de \unsure{le|la} faire)
On peut l'écrire

\[
\boxed{\deleted{\unsure{t}}T_{\varphi,\varphi'}\deleted{\unsure{(u,v)}}\unsure{(u,v)}=\mathrm{Tr}(v)\varphi(u)+\mathrm{Tr}(u)\varphi'(v)-2\deleted{\unsure{T}}\varphi(u)\varphi'(v).}
\]

Or, on a là une forme bil. en $u,v$,
qui pour $v=1$ se réduit :

\[
t(u,1)=\deleted{\unsure{\mathrm{Tr}}}2\varphi(u)+\mathrm{Tr}(u)-2(\varphi)=\mathrm{Tr}(u)
\]

et de \unsure{$=$|même} pour $u=1$ \deleted{\unsure{on}}

\[
t(u,1)=\mathrm{Tr}\,u,\qquad t(1,v)=\mathrm{Tr}\,v
\]

C'est donc un élément de $L_P$ ! Soit
\deleted{Je dis que $\unsure{\delta}$}
$\deleted{\unsure{t}}T_{\varphi,\varphi'}\in L_P$
défini par $\varphi,\varphi'$. J dis que

\[
\delta(t_{\varphi,\varphi'})=0
\]

Pour le voir, on peut faire le calcul
avec bases \deleted{\unsure{quelc.|quelconque}} $(1,u)$ $(1,v)$ telles que

\[
\varphi(u)=\varphi(v)=0.
\]

\inserted{d'où $\overset{c}{\overset{\parallel}{\det}}u=\overset{c'}{\overset{\parallel}{\det}}v=0$ car \deleted{\unsure{on prend}} $\varphi(u)=0$ \unsure{sol. de|de} $z^2-bz+c=0$}

de sorte que maintenant

\[
\deleted{\unsure{b'}}t(u,v)=0\deleted{\unsure{2\varphi}}.
\]

\deleted{Je dis que l'équation} Or on a l'identité (

\[
\begin{aligned}
\delta(u,v,t)&=t^2-bb't+(b^2c'+b'{}^2c-4cc')=\deleted{\unsure{t^2-bb't}}\\
&=t(t-bb')\qquad\text{car }c=c'=0\ \deleted{\text{car }c=c'=0}
\end{aligned}
\]

donc \deleted{en y faisant $t$} $t(u,v)=0$ en est
bien une solution.

